\documentclass[11pt]{article}

\usepackage{graphicx}
\usepackage[utf8]{inputenc}
\graphicspath{{./images/}}

\title{\includegraphics[scale=0.3]{logo.png} \\ \textbf{Data Link Protocol}}
\author{Computer Networks\\ Bachelors in Informatics and Computing Engineering \\ \\ 3LEIC03\_G3 \\ \\ Tiago Rodrigues up201907021@fe.up.pt \\ Mário Travassos up201905871@fe.up.pt  }
\date{\today}

\begin{document}

\maketitle

\newpage

\section*{Summary}

\paragraph{}This report will cover the first work proposed for the Computer Networks Curricular Unit, with the objective of creating a small application that could transfer data through two computers asynchronously, through a serial port.

\paragraph{}The application is capable of transferring files whilst mainting their integrity, and detect errors in transmission, resolving them if possible.

\section*{Introduction}

\paragraph{}This report is the result of an examination to the practical component, which was the development of a data transfer protocol. A serial port was used to transfer the files in an asynchronous fashion.

The report is organized as follows:

\begin{enumerate}
  \item{Architecture- Functional blocks and interfaces}
  \item{Code Structure - APIs, main code structures and their relation with the architecture}
  \item{Main use cases - Identification and Call Stack Sequence}
  \item{Data link Protocol - Main functional aspects and implementation strategy}
  \item{Application Protocol - Main functional aspects and implementation strategy }
  \item{Validation - Description of the tests conducted}
  \item{Efficiency - Statistical characterization of efficiency, against a Stop\&Wait protocol}
  \item{Conclusion - Summary of the above descriptions, reflection on the learning objectives}
\end{enumerate}

\section{Architecture}

The application consists of two main layers, one to interact with the file to be sent and another to interact with the hardware. They are the Application Layer, and the Data-Link layer, and they are detailed below.

\section{Code Structure}

\section{Main use cases}

\section{The Data link Layer}

\section{The Application Layer}

\section{Validation}

\section{Efficiency}

\section{Conclusion}

\end{document}
